% Input fragment. Requires hyperref (for \url), amsmath and amssymb.
% Companion bibliography entries are in references_extra.tex.

\section{Sources, prior methods, and the boundary of the contribution}
\label{sec:sources-literature}

\subsection{Positive interpolation is an established subject}

The construction in Section~\ref{sec:annihilation} does not originate
the general problem of Fourier interpolation by positive compactly
supported measures. Florek's 1987 paper has precisely that subject
in its title~\cite{Florek}. The publisher's bibliographic record was
available, but its full-text download returned an access error in
this investigation. Consequently, no specific quantitative theorem
is attributed to that paper and no claimed distinction rests on
guessing its contents.

As independently inspected context, Graham and Hare~\cite{GrahamHare}
study interpolation by nonnegative measures for Sidon,
$\varepsilon$-Kronecker and $I_0$ sets. Their introduction explicitly
identifies Florek's result as a Sidon-set predecessor. Such results
establish clear historical precedent for positive interpolation.
They do not, in the passages inspected, state the present support
bound for an arbitrary finite multiset in terms of its local
occupancy, with prescribed zero orders and a simultaneous uniform
tail estimate.

Kahane's treatment of nonharmonic Fourier series supplies a close
classical mechanism~\cite{Kahane}. Proposition~VI.2.1, on printed
page~131, gives a sufficient diagonal-dominance condition on the
Fourier Gram matrix of a compact positive measure for equivalence
between the coefficient $\ell^2$ norm and the associated integral
norm. Proposition~VI.2.2 gives the interpolation interpretation.
We use this as context for the matrix argument, without claiming
that those propositions themselves impose positivity on a new
interpolating density. In the present elementary proof the explicit
$\ell^1$ estimate on the correction coefficients is what preserves
positivity pointwise. Diagonal dominance, the Neumann series, and
convolution are all classical tools.

\subsection{Translate averaging and support bounds}

Kozma and Oravecz~\cite{KozmaOravecz}, Theorem~1, prove that a real
trigonometric polynomial with finite symmetric integer spectrum
$S\subset\mathbb Z^d\setminus\{0\}$ has a zero in every ball of
diameter $\sum_{\lambda\in S}(4\|\lambda\|_2)^{-1}$.
Their proof successively adds suitable translates to remove
frequencies while retaining positivity. The same elementary
frequency-removal operation is the dual counterpart of convolving
positive two-point measures. Although their stated theorem uses
integer frequencies, that removal step makes sense for real
frequencies as well. This is relevant precedent for a cost obtained
by summing reciprocal frequencies; it is not presented here as a
published local-occupancy bound for Dirichlet ordinates.

The comparison with the supplied arithmetic packet is narrower and
can be checked directly. Its explicit construction uses one
interval-convolution factor for each selected ordinate, with support
cost on the scale of $\sum_{0<\gamma\le B}\gamma^{-1}$, hence
$O((\log B)^2)$ for a fixed Dirichlet family. The construction proved
here groups occurrences into separated classes before convolution.
The local zero count bounds the number of classes by $O(\log B)$.
It thereby obtains support $O(\log B)$ while retaining exact
cancellation and adding the explicit tail attenuation. This is the
specific quantitative improvement over the inspected packet.

\subsection{The lower bound uses proved pair correlation only}

For multiplicity-preserving cancellation, Jensen's formula and the
Riemann--von Mangoldt count give the lower bound directly; simplicity
is irrelevant. For ordinary cancellation at distinct zeta ordinates,
Montgomery's proved squared-sinc pair-sum formula~\cite{Montgomery}
gives
\[
 \sum_{\gamma\ \mathrm{distinct},\ 0<\gamma\le T}m_\gamma^2
 \le (4/3+o(1))N(T).
\]
The parameter stays within the proved Fourier-support range and is
then allowed to approach its endpoint. No pair-correlation
conjecture is assumed. The integer inequality
$1\ge(3m-m^2)/2$ for $m\ge1$ yields the distinct-zero proportion
$5/6-o(1)$, also explicitly discussed by
Conrey--Ghosh--Gonek~\cite[pp.~497--498]{CGG}. This leads to the
convenient lower support constant $5/12$ for ordinary cancellation.
It is not asserted to be the strongest constant in the literature.

\subsection{Novelty and verification limits}

The claim supported by this packet is an explicitly proved
quantitative arithmetic improvement relative to the supplied
documents: logarithmic support, prescribed aggregate zero orders,
uniform high-frequency attenuation, and the resulting common
positive average of the smoothed residue-class errors. The
matching logarithmic obstruction explains the scale of the result.
The argument does not require independence of zero ordinates or
distinctness across different characters.

The literature search identified substantial related work and an
especially relevant older positive-interpolation paper whose full
text could not be inspected. Accordingly, global priority is
unverified. The packet claims neither a new general theory of
positive Fourier interpolation nor the solution of another named
conjecture. Numerical examples illustrate the finite construction;
the proofs, rather than numerical residuals, establish exact
cancellation.
